% Lecture section replacement command 
\newcounter{lectureCounter}
\renewcommand{\thesubsection}{\arabic{lectureCounter}.\arabic{subsection}}
\newcommand{\Lecture}[1]{
    \refstepcounter{lectureCounter}
    \setcounter{subsection}{0}
    \section*{Lecture \arabic{lectureCounter}: #1}
    \addcontentsline{toc}{section}{Lecture \arabic{lectureCounter}: #1}
}

% Theorem and definitions  
\newtheorem{theorem}{Theorem}[lectureCounter]
\newtheorem{lemma}[theorem]{Lemma}
\newtheorem{corollary}[theorem]{Corollary}
\newtheorem{definition}[theorem]{Definition}

\newenvironment{example}
  {\begin{proof}[Example]\renewcommand{\qedsymbol}{$\star$}}
  {\end{proof}}

\renewcommand{\thetheorem}{\arabic{lectureCounter}.\arabic{theorem}}

% Tikz environments 
\newenvironment{automata}{
    \begin{tikzpicture}[
        ->,
        >=Stealth,
        shorten >=1pt,
        auto,
        node distance=1.5cm,
        semithick,
        every state/.style={draw, minimum size=8mm}
    ]
}{
    \end{tikzpicture}
}

% Styling
\usepackage{xcolor}
\usepackage[
    colorlinks=true,
    linkcolor=black,
    urlcolor=blue!60!black,
]{hyperref}

% Macros imported from course website
\newcommand{\head}[1]{\begin{center}{\color{blue} \bf #1}\end{center}}

\newcommand{\ul}{\underline}
\newcommand{\ol}{\overline}

\newcommand{\defeq}{\stackrel {\mbox{\tiny def}}{=}}
\newcommand{\defequiv}{\stackrel {\mbox{\tiny def}}{\Longleftrightarrow}}
\newcommand{\set}[1]{\left\{ #1 \right\}}
\newcommand{\setdef}[2]{\left\{ #1 \mid #2 \right\}}
\newcommand{\just}[1]{\left\{ \mbox{#1} \right\}}
\newcommand{\tuple}[1]{\left< #1 \right>}


\newcommand{\langof}[1]{{\cal L}(#1)}
\newcommand{\length}[1]{\left| #1 \right|}
\newcommand{\size}[1]{\left| #1 \right|}
\newcommand{\numocc}[2]{\sharp #1 (#2)}
\newcommand{\eclass}[1]{\left[ #1 \right]}
\newcommand{\labeclass}[2]{\left[ #1 \right]_{#2}}
\newcommand{\stringop}[2]{\mbox{\bf #1 } #2}
\newcommand{\rev}[1]{\stringop{rev}{#1}}
\newcommand{\two}[1]{\stringop{two}{#1}}
\newcommand{\pref}[1]{\stringop{pref}{#1}}
\newcommand{\lpar}[1]{[_{#1}}
\newcommand{\rpar}[1]{]_{#1}}


\newcommand{\Not}[1]{\neg #1}
\newcommand{\Or}[2]{#1 \vee #2}
\newcommand{\And}[2]{#1 \wedge #2}
\newcommand{\Implies}[2]{#1 \Rightarrow #2}
\newcommand{\Equiv}[2]{#1 \Leftrightarrow #2}


\newcommand{\In}[2]{#1 \in #2}
\newcommand{\Notin}[2]{#1 \not\in #2}
\newcommand{\Union}[2]{#1 \cup #2}
\newcommand{\Intersect}[2]{#1 \cap #2}
\newcommand{\Product}[2]{#1 \times #2}
\newcommand{\Concat}[2]{#1 \cdot #2}
\newcommand{\Ast}[1]{{#1}^\ast}
\newcommand{\Astplus}[1]{{#1}^+}
\newcommand{\Pow}[1]{2^{#1}}
\newcommand{\Subset}[2]{#1 \subseteq #2}


\newcommand{\Mapping}[3]{#1 : #2 \rightarrow #3}
\newcommand{\refa}[2]{\alpha^{#1}_{#2}}
\newcommand{\subst}[3]{{\bf s}^{#1}_{#2} (#3)}
\newcommand{\Prod}[2]{#1 \rightarrow #2}
\newcommand{\OSDeriv}[2]{#1 \rightarrow_G #2}
\newcommand{\LabDeriv}[3]{#1 \stackrel{#2}{\rightarrow}_G #3}
\newcommand{\Deriv}[2]{\LabDeriv{#1}{\ast}{#2}}
\newcommand{\SDeriv}[2]{#1 \rightarrow_G #2}
\newcommand{\RDeriv}[2]{#1 \rightarrow^\ast_G #2}
\newcommand{\TDeriv}[2]{#1 \rightarrow^+_G #2}
\newcommand{\TDerivS}[3]{#1 \rightarrow^+_{#2} #3}
\newcommand{\Labprod}[3]{#1 \stackrel{#2}{\hookrightarrow} #3}
\newcommand{\Labtrans}[3]{#1 \stackrel{#2}{\longrightarrow} #3}
\newcommand{\Trans}[2]{#1 \rightarrow #2}
\newcommand{\RTtrans}[2]{#1 \rightarrow^{\ast} #2}
\newcommand{\config}[2]{\left< #1, #2 \right>}
\newcommand{\tmconfig}[3]{\left< #1, #2, #3 \right>}


